\section{Relation to prior work}
\label{sec:related}

\paragraph{Normalized persistence.}
Lowe et al.~\cite{lowe2026normalized} define normalized harmonic persistence
(their Problem~9, our \cref{def:problem}), prove that a low-energy relaxation
of it is \(\mathsf{DQC}_1\)-hard and in \(\BQP\), prove \(\SDQCone\)-hardness
of the local-Hamiltonian analogue, and conjecture \(\SDQCone\)-hardness of
Problem~9 (their Conjecture~1).  Their Conjecture~2 asks for a polynomial-time
map from history Hamiltonians \(H_1\) and \(H_2=H_1+\lambda H_{\rm out}\) to
weighted clique complexes \(X_1\subseteq X_2\) with two properties:
kernel fidelity and gap, \(\beta_d(X_i)=\dim\ker H_i\) together with
computable inverse-polynomial gap bounds; and functoriality,
\(\beta_d(X_1\to X_2)=\dim\ran(P_{\ker H_2}|_{\ker H_1})\).  Their Lemma~11
shows that Conjecture~2 implies Conjecture~1.  \Cref{thm:transfer,thm:quotient}
establish both properties for the exact \(G_2\) history Hamiltonians of
\cref{sec:circuit}, with \(H_{\rm out}\) at unit coefficient rather than at a
small coefficient \(\lambda\), and with the induced map computed through the
quotient model rather than through compatible harmonic isometries as
formulated there.  Two differences prevent us from concluding Conjecture~1
itself.  First, our source is \(\BQPone^{G_2}\) with three mixed label
qubits rather than an \(\SDQCone\) verifier with a maximally mixed register.
The \(\SDQCone\) promise constrains the acceptance probability of the
maximally mixed state, and a trace promise does not in general determine the
dimension of the perfectly accepted subspace, which is what \(q_d\) measures
(\cref{app:source-boundaries}).  Second, the certified palette is real and
exact, so verifiers over an arbitrary gate set are not covered.  Their
containment result (Lemma~7) needs additional access assumptions, namely a
preparation circuit for the uniform mixture over \(\ker\Delta_{1,d}\) and an
inverse-polynomial lower bound on the nonzero principal angles between the
two harmonic spaces.  We do not establish these for our instances, and
containment of the promise problem of \cref{def:problem} in \(\BQP\) remains
open.

\paragraph{Clique homology gadgets.}
Crichigno and Kohler~\cite{crichigno2024clique} give a parsimonious reduction
in which satisfying states correspond one-to-one to holes, so exact
multiplicity by itself is not new; their instances carry no gap promise.
King and Kohler~\cite{king2026gapped} introduce the weighted gadgets and the
many-gadget spectral analysis behind gapped clique homology, proving a gap
when the logical Hamiltonian is positive definite, with a gadget weight
proportional to the logical gap; Hayakawa~\cite{hayakawa2026unweighted}
removes the weights under the same positive-definite promise.  Our transfer
theorem is formulated from finite zero-weight certificates, applies to
arbitrary geometric chains and to degenerate kernels, and decouples
\(\lambda\) from \(g\).  Its padding step differs from the coordinate-bulk
formulation displayed in the arXiv version of King and
Kohler~\cite[Claim~10.2]{king2023gappedv2}: there the bulk consists of all
coordinates containing a private vertex, which in the padded complex also
contains boundary terms obtained by deleting a weight-one outside vertex,
whereas we keep local private coordinates tensored with outside harmonics
in every bidegree (\cref{lem:scaling-padding}).  We also work with the fixed
local matrices before tensoring with the outside identity, so that no
operator bound on a projector of growing rank has to be inferred from
basiswise perturbation estimates.  These are differences of proof technique,
and we make no claim about the presentation in the published
version~\cite{king2026gapped}.

\paragraph{Persistence hardness.}
Gyurik et al.~\cite{gyurik2026provable} prove \(\BQPone\)-hardness, over the
gate set \(\{\mathrm{CNOT},U_{\rm Pyth}\}\), of deciding whether a
\emph{given} harmonic representative of \(X_1\) survives in \(X_2\), with a
gap promise on the final Laplacian.  Normalized persistence asks instead for
the fraction of the full initial homology that survives, with no
representative given, so a hardness proof must control the complete kernels
at both endpoints.  Persistent Laplacians~\cite{memoli2022persistent} give a
spectral formulation of persistent Betti numbers; we use ordinary endpoint
Laplacians for the promise and the quotient model for the map.

\paragraph{Sources and gate sets.}
Rudolph~\cite{rudolph2026universal} defines the gate-dependent classes,
proves exact error reduction inside gate sets containing \(G_2\), shows
\(\BQPone^{G}\subseteq\BQPone^{G_2}\) for every finite \(G\) with entries in
some \(\mathbb Q(\zeta_{2^k})\), and supplies the gain/loss split of
\(H\otimes H\) with rational three-term constraint vectors and the base
clique gadgets used in our palette.  Cade and
Crichigno~\cite{cade2024complexity} relate cohomology problems to
\(\QMAone\) and \(\mathsf{DQC}_1\), and Gyurik, Cade, and
Dunjko~\cite{gyurik2022towards} connected quantum topological data analysis
to \(\mathsf{DQC}_1\)~\cite{knill1998power}.  The history-state construction
with a unary clock and neighbor guards is
Kitaev's~\cite{kitaev2002classical,kempe2006complexity}, with the rational
amplitude bookkeeping of \cref{sec:circuit}.

\paragraph{Unweighting.}
Hayakawa's symmetric/asymmetric decomposition of the common-copy
blow-up~\cite[Lemma~4.2 and Theorem~4.3]{hayakawa2026unweighted} is used as a
black box.  Our only addition is the observation that identical labeled copy
blocks at both filtration levels make the symmetric isometries commute with
inclusion, \cref{eq:unweight-natural}.  We do not count this as a separate
contribution.
